\documentclass[14pt,a4paper,oneside]{extbook}
\usepackage[utf8]{vietnam}
\usepackage{amsmath,amssymb,mathrsfs,fancyhdr,enumerate,multirow,makecell}
\usepackage{tkz-euclide,tikz-3dplot,tikz,tkz-tab}
\usetikzlibrary{arrows.meta,calc,intersections,angles,quotes,backgrounds,shapes.geometric,patterns,through,shadows,positioning}
\usepackage{yhmath}
\usepackage[top=0.7cm, bottom=0.7cm, left=1.0cm, right=1.0cm, headheight=16pt, headsep=3mm, footskip=5mm]{geometry}
\usepackage[hidelinks,unicode]{hyperref}
\usepackage{qrcode}
\usepackage[many]{tcolorbox}
\usepackage{fontawesome}

\pagestyle{fancy}
\fancyhf{}
\fancyhead[L]{\footnotesize\bfseries\color{teal!85!black}\faBook~SỔ TAY KIẾN THỨC TOÁN 5}
\fancyhead[R]{\footnotesize\bfseries\color{blue!85!black}Lớp Toán Thầy Lĩnh -- loptoanthaylinh.vn}
\fancyfoot[C]{\footnotesize\bfseries\color{gray!80!black}-- Sổ tay kiến thức: Bài 63. Thu thập, phân loại, sắp xếp số liệu --}
\renewcommand{\headrulewidth}{0.4pt}

\begin{document}

% --- BANNER TIÊU ĐỀ BÀI HỌC ---
\noindent
\begin{tcolorbox}[
  enhanced,
  colback=teal!8!white,
  colframe=teal!80!black,
  arc=3mm,
  boxrule=1.1pt,
  width=\linewidth,
  halign=center,
  top=2.2mm, bottom=2.2mm,
  left=3mm, right=3mm,
  drop fuzzy shadow
]
  {\large\bfseries\color{teal!95!black}\faBookmark\ BÀI 63. THU THẬP, PHÂN LOẠI, SẮP XẾP SỐ LIỆU}\\[2pt]
  {\scriptsize\color{black!75}\itshape Quy trình xử lý số liệu thống kê, kĩ năng lập bảng và kiểm đếm trên 1 trang A4}
\end{tcolorbox}

\vspace{3.2mm}
% --- KHỐI 1: QUY TRÌNH XỬ LÝ SỐ LIỆU ---
\noindent
\begin{minipage}[t]{0.56\textwidth}
  \begin{tcolorbox}[
    enhanced,
    title={\footnotesize\bfseries\color{white}\faListOl\ 1. Quy trình 3 bước xử lý số liệu},
    colback=blue!2!white,
    colframe=blue!75!black,
    arc=2.5mm, boxrule=0.8pt, top=1.8mm, bottom=1.8mm, left=2.2mm, right=2.2mm
  ]
    \footnotesize
    $\bullet$ \textbf{Bước 1 (Thu thập):} Quan sát, đo đạc, phỏng vấn hoặc kiểm đếm để lấy số liệu ban đầu.\\
    $\bullet$ \textbf{Bước 2 (Phân loại):} Chia các số liệu theo tiêu chí định trước (màu sắc, giới tính, kích cỡ...).\\
    $\bullet$ \textbf{Bước 3 (Sắp xếp):} Xếp theo thứ tự (tăng dần, giảm dần) và đưa vào bảng biểu chuẩn.
  \end{tcolorbox}
\end{minipage}%
\hfill
\begin{minipage}[t]{0.42\textwidth}
  \begin{tcolorbox}[
    enhanced,
    title={\footnotesize\bfseries\color{white}\faTable\ 2. Kĩ năng đọc bảng số liệu},
    colback=teal!2!white,
    colframe=teal!75!black,
    arc=2.5mm, boxrule=0.8pt, top=1.8mm, bottom=1.8mm, left=2mm, right=2mm
  ]
    \footnotesize
    $\bullet$ \textbf{Tìm giá trị cực trị:} Đối tượng có số lượng nhiều nhất, ít nhất.\\
    $\bullet$ \textbf{So sánh chênh lệch:} Nhiều hơn/kém hơn bao nhiêu đơn vị, gấp mấy lần.\\
    $\bullet$ \textbf{Tính tổng số đối tượng:} Cộng tổng số liệu của tất cả các nhóm.
  \end{tcolorbox}
\end{minipage}

\vspace{3mm}

% --- KHỐI 3: SƠ ĐỒ TƯ DUY ---
\noindent
\begin{tcolorbox}[
  enhanced,
  title={\footnotesize\bfseries\color{white}\faSitemap\ Sơ đồ tư duy: Quy trình xử lý số liệu thống kê},
  colback=yellow!2!white,
  colframe=teal!85!black,
  arc=3mm, boxrule=0.9pt, top=1.8mm, bottom=1.8mm, left=1.5mm, right=1.5mm,
  halign=center
]
  \begin{tikzpicture}[
    scale=0.92,
    every node/.style={transform shape},
    b0/.style={rectangle, rounded corners=5pt, fill=teal!85!black, text=white, font=\bfseries\footnotesize, inner sep=5pt, align=center},
    b1/.style={rectangle, rounded corners=4pt, fill=blue!15, draw=blue!70!black, text=blue!90!black, font=\bfseries\scriptsize, inner sep=4pt, align=center},
    b2/.style={rectangle, rounded corners=4pt, fill=orange!15, draw=orange!70!black, text=orange!90!black, font=\bfseries\scriptsize, inner sep=4pt, align=center},
    b3/.style={rectangle, rounded corners=4pt, fill=green!15, draw=green!70!black, text=green!90!black, font=\bfseries\scriptsize, inner sep=4pt, align=center}
  ]
    \node[b0] (core) at (-6.2, 0) {\faDatabase\ THỐNG KÊ\\SỐ LIỆU};
    \node[b1] (n1) at (-2.0, 0) {\textbf{BƯỚC 1: THU THẬP}\\[2pt]Quan sát, phỏng vấn,\\kiểm đếm số liệu};
    \node[b2] (n2) at (2.2, 0) {\textbf{BƯỚC 2: PHÂN LOẠI}\\[2pt]Chia nhóm theo\\tiêu chí rõ ràng};
    \node[b3] (n3) at (6.2, 0) {\textbf{BƯỚC 3: SẮP XẾP}\\[2pt]Lập bảng số liệu\\(Tăng/giảm dần)};

    \draw[teal!85!black, thick, -{Stealth[length=2.5mm]}] (core) -- (n1);
    \draw[teal!85!black, thick, -{Stealth[length=2.5mm]}] (n1) -- (n2);
    \draw[teal!85!black, thick, -{Stealth[length=2.5mm]}] (n2) -- (n3);
  \end{tikzpicture}
\end{tcolorbox}

\vspace{3mm}

% --- KHỐI 4: PHƯƠNG PHÁP KIỂM ĐẾM & BẪY TRẮC NGHIỆM ---
\noindent
\begin{minipage}[t]{0.54\textwidth}
  \begin{tcolorbox}[
    enhanced,
    title={\footnotesize\bfseries\color{white}\faCheckSquareO\ Phương pháp gạch sổ cụm 5 kiểm đếm},
    colback=green!3!white,
    colframe=green!60!black,
    arc=2.5mm, boxrule=0.8pt, top=1.8mm, bottom=1.8mm, left=2.2mm, right=2.2mm
  ]
    \footnotesize
    $\bullet$ Gạch sổ theo cụm $5$ gạch: $4$ gạch đứng và $1$ gạch chéo (\textsf{/\!/\!/\!/}\hspace{-6pt}\textbf{---}) giúp đếm cực nhanh không sót.\\
    $\bullet$ Khi đếm xong, lấy số cụm nhân $5$ rồi cộng số gạch lẻ còn lại.
  \end{tcolorbox}
\end{minipage}%
\hfill
\begin{minipage}[t]{0.44\textwidth}
  \begin{tcolorbox}[
    enhanced,
    title={\footnotesize\bfseries\color{white}\faExclamationTriangle\ Mẹo tránh bẫy trắc nghiệm},
    colback=red!2!white,
    colframe=red!70!black,
    arc=2.5mm, boxrule=0.8pt, top=1.8mm, bottom=1.8mm, left=2.2mm, right=2.2mm
  ]
    \footnotesize
    $\bullet$ \textbf{Bẫy thiếu/thừa đối tượng:} Luôn cộng tổng số lượng các nhóm phân loại xem có bằng đúng tổng số đối tượng điều tra ban đầu không!\\
    $\bullet$ \textbf{Bẫy tiêu chí:} Mỗi đối tượng chỉ được thuộc đúng $1$ nhóm duy nhất.
  \end{tcolorbox}
\end{minipage}

\end{document}
